% ============================================================================
% EXAMPLE: Single-paper conference poster (auto-generated)
%
% Produced by running:
%   python scripts/plan_poster.py /tmp/bundle_test \
%     --output-dir /tmp/test_poster \
%     --title "A retrieved-context account of contiguity effects" \
%     --size 48x36-landscape
%
% on a fake paper bundle. Preserved here as a reference for what the skill
% emits before deep-reading fills in the [TODO ...] placeholders. Compiles
% to a 48x36 landscape PDF with pdflatex (no bibtex needed — refs are
% hand-formatted inline).
% ============================================================================
% ============================================================================
% CML Beamerposter Template
% Matches the visual style of CML Beamer decks (Madrid theme + navy blue).
% Designed for standard 48x36 landscape conference posters.
%
% Variables to fill in (look for <<...>> markers):
%   A retrieved-context account of contiguity effects          main poster title (can be 1 or 2 lines)
%   Alice Smith$^1$, Bob Jones$^1$, Carol Kim$^1$        presenter names with superscript affiliation markers
%   $^1$University of Pennsylvania   matching superscript-numbered affiliations
%             conference or meeting name (optional; can be blank)
%
% Change poster size by editing the beamerposter \usepackage options.
% Size presets:
%   48x36 landscape (US default): width=121.9, height=91.4
%   36x48 portrait:               width=91.4,  height=121.9
%   A0 portrait:                  size=a0, orientation=portrait (no width/height)
%   A0 landscape:                 size=a0, orientation=landscape
%
% Column count: this template assumes 3 columns. For 2 columns change all
% \begin{column}{0.32\textwidth} to {0.48\textwidth}.
% ============================================================================

\documentclass[final]{beamer}
\usepackage[orientation=landscape,size=custom,width=121.9,height=91.4,scale=1.35]{beamerposter}

% --- Theme ------------------------------------------------------------------
\usetheme{Madrid}
\usecolortheme{default}

% CML navy blue — same shade as the Beamer template
\definecolor{cmlblue}{RGB}{28, 53, 115}
\setbeamercolor{structure}{fg=cmlblue}
\setbeamercolor{palette primary}{bg=cmlblue, fg=white}
\setbeamercolor{palette secondary}{bg=cmlblue!80, fg=white}
\setbeamercolor{palette tertiary}{bg=cmlblue!60, fg=white}

% Make block titles more prominent on a poster than they are on slides
\setbeamercolor{block title}{fg=white, bg=cmlblue}
\setbeamercolor{block body}{fg=black, bg=cmlblue!5}
\setbeamerfont{block title}{series=\bfseries, size=\Large}

% beamerposter with Madrid already hides navigation chrome — don't override
% \setbeamertemplate{navigation symbols}{}, \footline, or \headline here.
% Those overrides confuse beamerposter's page-geometry calculation and
% trigger "Illegal unit of measure" warnings at \end{frame}.

% --- Packages ---------------------------------------------------------------
\usepackage[T1]{fontenc}
\usepackage[utf8]{inputenc}
\usepackage{helvet}  % scalable Helvetica-compatible sans — the conventional poster font
\renewcommand{\familydefault}{\sfdefault}  % use sans as body default
\usepackage{amsmath, amssymb}
\usepackage{graphicx}
\usepackage{booktabs}
\usepackage{tcolorbox}
\tcbuselibrary{skins}
% \usepackage{qrcode}  % uncomment if including a QR code

% Spacing in blocks — posters breathe better with a bit more gap
\setlength{\parskip}{0.5em}
\setlength{\itemsep}{0.3em}

% ============================================================================
\begin{document}

\begin{frame}[t]{}

% --- Title header (full-width banner) ---------------------------------------
\begin{tcolorbox}[colback=cmlblue, colframe=cmlblue, arc=5pt,
                  boxsep=12pt, left=20pt, right=20pt, top=8pt, bottom=8pt]
  \begin{center}
    {\fontsize{80}{96}\selectfont\bfseries\color{white} A retrieved-context account of contiguity effects} \\[0.5em]
    {\Large\color{white} Alice Smith$^1$, Bob Jones$^1$, Carol Kim$^1$} \\[0.3em]
    {\normalsize\color{white} $^1$University of Pennsylvania}
  \end{center}
\end{tcolorbox}
\vspace{1.5em}

% --- Three columns ---------------------------------------------------------
\begin{columns}[T]

% Generated by scripts/plan_poster.py from the poster plan.

\begin{column}{0.32\textwidth}

\begin{block}{Motivation}
  [TODO — 40-80 words. Lead with the phenomenon or puzzle, not the literature.]
\end{block}



\begin{block}{Question}
  \begin{center}
    \bfseries [TODO — one-sentence research question, visually emphasized.]
  \end{center}
\end{block}



\begin{block}{Methods}
  [TODO — 50-100 words. Task, N, key DV. Prefer a schematic to prose.]
\end{block}



\end{column}

\begin{column}{0.32\textwidth}

\begin{block}{[TODO — takeaway-style heading for fig1.pdf]}
  \begin{center}
    \includegraphics[width=0.95\textwidth]{figures/fig1.pdf}
  \end{center}
  \vspace{0.3em}
  {\small [TODO — 30-80 words. What to notice and what it means.]}
\end{block}



\begin{block}{[TODO — takeaway-style heading for fig2.pdf]}
  \begin{center}
    \includegraphics[width=0.95\textwidth]{figures/fig2.pdf}
  \end{center}
  \vspace{0.3em}
  {\small [TODO — 30-80 words. What to notice and what it means.]}
\end{block}



\end{column}

\begin{column}{0.32\textwidth}

\begin{block}{[TODO — takeaway-style heading for fig3.pdf]}
  \begin{center}
    \includegraphics[width=0.95\textwidth]{figures/fig3.pdf}
  \end{center}
  \vspace{0.3em}
  {\small [TODO — 30-80 words. What to notice and what it means.]}
\end{block}



\begin{block}{[TODO — takeaway-style heading for fig4.pdf]}
  \begin{center}
    \includegraphics[width=0.95\textwidth]{figures/fig4.pdf}
  \end{center}
  \vspace{0.3em}
  {\small [TODO — 30-80 words. What to notice and what it means.]}
\end{block}



\begin{block}{Conclusions}
  \begin{itemize}
    \item {[TODO — complete-sentence claim 1.]}
    \item {[TODO — complete-sentence claim 2.]}
    \item {[TODO — complete-sentence claim 3.]}
  \end{itemize}
\end{block}



\begin{block}{References \& Acknowledgments}
  {\scriptsize
    \textbf{References:}\par
    {}[1] [TODO — reference 1]\par
    {}[2] [TODO — reference 2]\par
    \vspace{0.5em}
    \textbf{Acknowledgments:} [TODO — grant numbers, collaborators, contact.]
  }
\end{block}



\end{column}

\end{columns}

\end{frame}
\end{document}
